Sparsification would still be worth having if it cost nothing in return. Testing
that requires care, because on a task where the channel is not needed
``sparsification is free'' and ``the channel was never used'' make identical
predictions. We therefore gate the study on a positive control and register the
analysis plan before the outcomes are inspected. The study is secondary. Its
role is to establish that the architecture being sparsified does something, so
that the sparsification comparisons are not vacuous, and to report what a
five-seed design can and cannot resolve.

\subsection{Design}

Twelve variants at up to twelve agents, five seeds each, two million
environment steps. All scaling in the number of agents is delegated to the
training-free measurements of Section~\ref{sec:measurement}, because an
end-to-end training step in this pipeline is dominated by the environment
(Section~\ref{sec:backward}) and cannot attribute a latency difference to the
communication module.

Two positive controls establish that the channel carries something. The first is
a reference task in which each agent observes only its partner's goal, so
information must pass between agents; it is fixed at two agents, where the
selection mechanism is degenerate, so it validates the message pathway rather
than the selection. The second is a partially observable variant of a standard
cooperative task, introduced here, in which a peer's relative position is
withheld beyond a fixed radius. Under a uniform random policy the fraction of
peers visible falls with the number of agents while the absolute number visible
rises, because the world does not grow: that is a difficulty confound across
agent counts, and it is why no claim in this paper compares returns across agent
counts in this environment.

The analysis plan was committed after eight of the forty runs in the
pre-registered family had completed and before any of their outcomes was
inspected; the two positive controls were complete and had been inspected, which
is why they are excluded from the inferential family.
Appendix~\ref{app:prereg} gives the chronology and the single amendment. The
plan fixes the endpoint, the six comparisons and their direction, an exact
two-sided permutation test, the effect sizes, a percentile bootstrap interval
flagged approximate at five seeds, and a Holm--Bonferroni correction. It also
fixes the power expectation in advance: with five seeds per variant the smallest
attainable two-sided $p$-value is $2/252 = 0.0079$ against a first Holm
threshold of $0.05/6 = 0.0083$, a margin of $4\times10^{-4}$, so a first
rejection requires one variant's five seeds to lie entirely above the other's.
Any null in this family is under-powered rather than evidence of absence, and no
equivalence claim is available because no equivalence margin was registered.

\subsection{Results}

\begin{table}[t]
\centering\small
\caption{The two positive controls, which were inspected before registration and
are excluded from the correction, and the primary pre-registered comparison.
$A-B$ is the difference of means and $d$ the standardised effect.
Appendix~\ref{app:training} gives the remaining five comparisons.}
\label{tab:training}
\begin{tabular}{lrrrr}
\toprule
& $A-B$ & 95\% interval & $d$ & $p$ \\
\midrule
\multicolumn{5}{l}{\emph{Positive controls, descriptive}} \\
communicating $-$ non-communicating, reference task
  & $+4.26$ & $[+3.15, +5.23]$ & --- & $0.0079$ \\
communicating $-$ non-communicating, partially observable
  & $+11.59$ & $[+7.13, +16.44]$ & --- & $0.0079$ \\
\midrule
\multicolumn{5}{l}{\emph{Primary pre-registered comparison}} \\
attention-selected peers $-$ uniformly random peers, same budget
  & $+4.03$ & $[-3.35, +11.42]$ & $+0.60$ & $0.357$ \\
\bottomrule
\end{tabular}
\end{table}

Both positive controls separate completely, each attaining the floor of the
permutation test for this design (Table~\ref{tab:training}). Measured attention
entropy on the partially observable control is $0.35$ nats against a ceiling of
$1.61$, so the learned weighting is peaked rather than uniform. What this
establishes is that sparsification is being applied to an architecture whose
communication machinery improves return under the evaluated setup. It does not
establish that the learned attention channel alone accounts for the difference:
the communicating variants carry $7$--$19\%$ more trainable parameters, and the
reference task supplies a discrete utterance channel of its own. A
parameter-matched sham-communication variant would be needed for the
channel-specific causal claim, and we do not make it.

The comparison the design exists for holds the communication budget, the
executed aggregation cost and every hyperparameter fixed, and varies only
whether peers are chosen by attention score or uniformly at random. Its point
estimate favours attention by roughly two thirds of a standard deviation and its
interval includes zero. The design lacked the power to resolve an effect of that
size, which the pre-registered power calculation said in advance; nothing
follows about whether attention selects informatively. The remaining five
comparisons are reported in Appendix~\ref{app:training} as point estimates with
intervals \citep{agarwal2021rliable}, and all six are Holm-adjusted nulls. No
conclusion in this paper rests on any of them.

Two limits of the environment bound what the study can say. At twelve agents
neither variant learns the task well: the dense variant scores below a
uniform-random policy and the sparse variant only just above it, so no
comparison at that agent count is interpretable. And the centralised critic sees
the same masked observations as the actors rather than a privileged global
state, which departs from standard practice for this learner and is a plausible
contributor to that failure.
